% !TEX root = ../../../main.tex
% Appendix entry: preprocessing.detrend -- "Baseline removal" (Appendix A).
\entryfield{Algorithm}
\code{_run} in \code{Adapters/preprocessing_detrend.py} calls the factory
\code{make_detrend_filter} in \code{Working/Preprocessing/detrend.py} with the span's $f_s$, the
\emph{mode} and \emph{window\_s}, and applies the function it returns. The input is cast to
float64. Let $W = \max(1, \mathrm{round}(\mathit{window\_s}\cdot f_s))$ samples. By mode:
\begin{description}
\item[\texttt{rolling\_mean}] With $h = \max(1, \lfloor W/2 \rfloor)$, the trend at sample $i$
  is the mean of the finite samples in $[\max(0, i-h),\ \min(n-1, i+h)]$ (function
  \code{_rolling_mean}, via \code{_finite_window_mean}, from cumulative sums of the finite
  values and of their count). The window holds $2h+1$ samples (601 at the defaults on a 1\,Hz
  recording) and shrinks to what fits near the ends, becoming one-sided there. The output is
  $y_i = x_i - m_i$.
\item[\texttt{rolling\_mean\_nearest}] The drop detector's filter, imported rather than
  re-implemented: \code{detrend} in \code{Working/Detection/drop_motifs/detect.py}. The window
  is $\max(3, W)$ samples. On finite input it is $x -$
  \code{scipy.ndimage.uniform_filter1d(x, size, mode="nearest")}: the ends are padded by
  repeating the end sample, so every window has full width; for an even size SciPy's window
  covers $i - \lfloor\mathrm{size}/2\rfloor$ to $i + \lfloor\mathrm{size}/2\rfloor - 1$. When
  the window is at least the span's length, the span's mean is subtracted instead. With NaNs, the
  same padding and alignment are used but each window averages only its finite samples.
\item[\texttt{rolling\_z}] First $r_i = x_i - m_i$ as in \texttt{rolling\_mean}; then
  $s_i = \sqrt{\overline{r^2}_i}$, the same rolling mean applied to $r^2$; the output is
  $r_i / s_i$ where $s_i > 10^{-12}$ and 0 elsewhere (a NaN residual stays NaN). Note that
  $s_i$ is the root mean square of residuals measured about \emph{each sample's own} trend,
  which approximates, but is not exactly, the rolling standard deviation.
\item[\texttt{linear}] One least-squares line over the whole span,
  \code{np.polyfit(idx, x, 1)}, subtracted; \emph{window\_s} is ignored.
\end{description}
Any other mode raises a \code{ValueError} when the filter is built.

\entryfield{Assumptions}
\begin{itemize}
\item \emph{Sampling rate.} \emph{window\_s} is in seconds and converted with the span's
  $f_s$. The registry's minimum is 1\,s.
\item \emph{Units.} Not unit-aware. The rolling-mean and linear modes return the stored
  unit; \texttt{rolling\_z} returns a unit-free z-score, so a threshold downstream means
  standard deviations, not volts.
\item \emph{Missing values.} In the three rolling modes a NaN (or infinite) sample stays NaN,
  the samples beside it are detrended against the finite samples in their window, and the
  output is NaN only where a window holds no finite sample, so a gap keeps its own length
  (the 2026-10-14 fix; before it, one NaN blanked every later window). In \texttt{linear} mode
  NaNs reach \code{np.polyfit}, which the docstring says raises.
  \unverified{The exact error has not been checked.}
\item \emph{Span length.} A window longer than the span is allowed: \texttt{rolling\_mean}
  then subtracts the span's mean everywhere (every window covers the whole span), as does
  \texttt{rolling\_mean\_nearest} explicitly.
\end{itemize}

\entryfield{Outputs}
A \texttt{Signal} of the input's length and $f_s$; nothing in \code{meta}.

\entryfield{Cost}
$O(n)$ time and memory in every mode: the rolling modes use cumulative sums (or SciPy's running
filter), independent of the window length. No estimate is registered.

\entryfield{Failure modes}
\begin{itemize}
\item \emph{Overshoot beside large events.} The moving average includes the events themselves:
  a deep drop pulls the local mean down, so after subtraction the trace on either side of the
  drop is lifted and the drop looks shallower. \code{tests/test_beads.py} measures this on its
  synthetic acceptance trace (drops of depth 1 recovering with $\tau = 200$\,s, a 600\,s
  window): an excursion of about 0.34 of the drop depth at every drop, against about 0.06 for
  \blockref{preprocessing.beads}.
\item \emph{Events longer than the window are cut down.} A slow recovery longer than about
  half the window is treated as trend. The literature notes record this as a repeated finding
  on two event families (\code{docs/research/analysis-blocks-from-literature.md}).
\item \emph{Ends (\texttt{rolling\_mean}).} The shrinking, one-sided window near each end
  follows the data more closely there, so the first and last half-window are flattened more
  than the interior. \texttt{rolling\_mean\_nearest} instead repeats the end sample, which
  avoids an artificial step at the edges (the reason the drop detector uses it).
\item \emph{\texttt{rolling\_z} amplifies quiet stretches.} Where the signal is nearly flat the
  local spread is tiny and the z-score inflates the noise to unit size; a perfectly flat window
  gives 0.
\item \emph{\texttt{linear} on curved drift} leaves the curvature in place.
\end{itemize}

\entryfield{Validation}
\code{tests/test_detrend_nan.py} checks that finite input gives exactly the numbers of the
earlier, non-NaN-aware implementation (kept verbatim in the test), that each rolling mode's
output equals a one-window-at-a-time reference, and the NaN behaviour above; it also checks
that \texttt{rolling\_mean\_nearest} is exactly the drop detector's filter.
\code{tests/test_template_drop_detection.py} checks that the drop-detection template, which
starts with this block in \texttt{rolling\_mean\_nearest} mode, reproduces the seed events. The
adapter's docstring cites a Stage 3 benchmark (SAX recall 0.22 raw against 0.85 detrended);
the evidence behind that number was not traced here.

\entryfield{References}
None: moving-average and least-squares detrending are standard. Why the project prefers BEADS
for event-dominated channels is set out under \blockentry{preprocessing.beads}.

\entryfield{Tests}
\code{tests/test_detrend_nan.py} (the filter itself), \code{tests/test_t06_adapter_remap.py}
(types, length, rate), \code{tests/test_drop_detection_nan.py} and the
\code{tests/test_template_*.py} files that start their chains with this block,
\code{tests/test_beads.py} (the comparison above) and \code{tests/test_webui_block_readout.py}.
